\chapter{Stima bayesiana: MMSE e MAP}

Nei due capitoli precedenti il parametro da scoprire poteva assumere pochi valori, uno per ipotesi, e il problema era
\textbf{decidere}. Spesso invece il parametro è una grandezza continua (la probabilità che una moneta dia testa, la
distanza di un bersaglio, la temperatura di un forno) e il problema è \textbf{stimarlo}: dare un numero il più vicino
possibile a quello vero. Ci sono due approcci. In quello \textbf{bayesiano}, di questo capitolo, il parametro è a sua
volta una variabile aleatoria, di cui si conosce la densità a priori. In quello \textbf{classico}, del prossimo
capitolo, è una quantità non nota ma deterministica.

% ============================================================
\section{Il problema della stima}
% ============================================================

Si abbia un insieme di dati $\vett{x} \in \mathcal{X}^n$, realizzazione di un vettore aleatorio $\vett{X}$ la cui
distribuzione appartiene a una famiglia di pmf $p_{\vett{X} \mid \Theta}(\vett{x} \mid \theta)$ o di PDF
$f_{\vett{X} \mid \Theta}(\vett{x} \mid \theta)$, dove il valore di $\theta$ non è noto. Tipicamente $\theta$ è un
parametro continuo, che può essere
\begin{itemize}
  \item la realizzazione di una variabile aleatoria continua $\Theta$, con una densità marginale $f_\Theta(\theta)$ nota
        (\textbf{approccio bayesiano}): la densità \term{a priori}, che descrive che cosa si sa del parametro prima di
        guardare i dati;
  \item una quantità non nota e deterministica, che assume valori in un insieme continuo (\textbf{approccio classico}).
\end{itemize}
Domanda: come si stima $\theta$ sulla base del campione osservato?

Nell'approccio bayesiano i dati aggiornano quello che si sa del parametro. Per la legge di Bayes (con le densità del
Capitolo~14), la densità \term{a posteriori} del parametro, sapendo i dati, è
\[ f_{\Theta \mid \vett{X}}(\theta \mid \vett{x}) = \begin{cases} \dfrac{p_{\vett{X} \mid \Theta}(\vett{x} \mid
   \theta)\, f_\Theta(\theta)}{\int p_{\vett{X} \mid \Theta}(\vett{x} \mid \theta)\, f_\Theta(\theta) \dd\theta} &
   \text{dati discreti,} \\[5mm] \dfrac{f_{\vett{X} \mid \Theta}(\vett{x} \mid \theta)\, f_\Theta(\theta)}{\int
   f_{\vett{X} \mid \Theta}(\vett{x} \mid \theta)\, f_\Theta(\theta) \dd\theta} & \text{dati continui,} \end{cases} \]
dove i denominatori sono la legge non condizionata dei dati, per la probabilità totale:
\[ p_{\vett{X}}(\vett{x}) = \int p_{\vett{X} \mid \Theta}(\vett{x} \mid \theta)\, f_\Theta(\theta) \dd\theta , \qquad
   f_{\vett{X}}(\vett{x}) = \int f_{\vett{X} \mid \Theta}(\vett{x} \mid \theta)\, f_\Theta(\theta) \dd\theta . \]
(Se $\Theta$ fosse una variabile discreta, si ricadrebbe in un problema di classificazione come quelli del
Capitolo~19.)

\dfn{Stimatore}{
  Uno \term{stimatore} del parametro $\theta$ è una variabile aleatoria $\hat\Theta(\vett{X})$, funzione dei dati, le
  cui realizzazioni $\hat\theta(\vett{x})$ provano a «indovinare» il valore di $\theta$ sulla base di un'osservazione
  $\vett{x} \in \mathcal{X}^n$. Il numero $\hat\theta(\vett{x})$ si chiama \term{stima}.
}

Il cappello $\hat{\ }$ si legge «stimato»: $\hat\theta$ è «theta stimato». Per progettare uno stimatore bisogna dire
quanto costano gli errori, come per le decisioni.

\dfn{Rischio bayesiano di uno stimatore}{
  Data una \term{funzione di costo} $C(\cdot)$ dell'errore di stima, il \term{rischio medio bayesiano} dello stimatore
  $\hat\Theta$ è il costo medio:
  \[ R = E\left[ C\left( \hat\Theta(\vett{X}) - \Theta \right) \right] = E_{\vett{X}}\left[ E\left[ C\left(
     \hat\Theta(\vett{X}) - \Theta \right) \Bigm| \vett{X} \right] \right] . \]
  Uno \term{stimatore ottimo} è uno stimatore che minimizza il rischio bayesiano.
}

La seconda forma del rischio è il teorema della media condizionata (Capitoli~9 e~14): prima si fa la media sul
parametro sapendo i dati, poi sui dati.

\begin{richiamo}{arg min e arg max}
Per una funzione $g(t)$, $\min_t g(t)$ è il \textbf{valore} minimo della funzione, mentre $\arg\min_t g(t)$
(«l'argomento del minimo») è il \textbf{punto} $t$ in cui il minimo viene raggiunto. Per esempio, per $g(t) = (t -
3)^2 + 5$ si ha $\min_t g(t) = 5$ ma $\arg\min_t g(t) = 3$. Allo stesso modo $\arg\max$ è il punto di massimo. Due
osservazioni utili: moltiplicare $g$ per una costante positiva o sommarle una costante non sposta il punto di minimo
(cambia solo il valore); e applicare a $g$ una funzione strettamente crescente, come il logaritmo, non sposta né il
minimo né il massimo.
\end{richiamo}

\teorema[stimatoreottimo]{Lo stimatore bayesiano ottimo}{
  Lo stimatore ottimo bayesiano che opera sul campione osservato $\vett{x}$ è
  \[ \hat\theta(\vett{x})_{\text{opt}} = \arg\min_{\hat\theta} \int C\left( \hat\theta - \theta \right) f_{\Theta \mid
     \vett{X}}(\theta \mid \vett{x}) \dd\theta : \]
  per ogni $\vett{x}$ si sceglie il valore che rende minimo il costo medio \emph{a posteriori}.
}

\begin{dimostrazione}
\passo{Il rischio come somma sui dati} Con dati discreti, scrivendo la media esterna sui dati:
\[ E\left[ C\left( \hat\Theta(\vett{X}) - \Theta \right) \right] = \sum_{\vett{x} \in \mathcal{X}^n}
   p_{\vett{X}}(\vett{x}) \int C\left( \hat\theta(\vett{x}) - \theta \right) f_{\Theta \mid \vett{X}}(\theta \mid
   \vett{x}) \dd\theta . \]
\passo{Si minimizza un addendo alla volta} La stima $\hat\theta(\vett{x})$ si può scegliere separatamente per ogni
$\vett{x}$, e compare solo nel proprio addendo. Se per ogni $\vett{x}$ si prende il minimo dell'integrale, pesando con
le quantità non negative $p_{\vett{X}}(\vett{x})$ e sommando si ottiene il minimo globale. (Con dati continui la somma è
un integrale in $\dd\vett{x}$, e il ragionamento è lo stesso.)
\end{dimostrazione}

% ============================================================
\section{Costo quadratico: lo stimatore MMSE}
% ============================================================

La scelta più comune è un costo \textbf{quadratico}, $C\left( \hat\Theta(\vett{X}) - \Theta \right) = \left(
\hat\Theta(\vett{X}) - \Theta \right)^2$: gli errori piccoli costano poco, quelli grandi moltissimo. Il rischio è
allora l'errore quadratico medio, e lo stimatore ottimo si chiama \term{stimatore MMSE} (\emph{Minimum Mean Square
Error}, a minimo errore quadratico medio).

\begin{richiamo}{derivare sotto il segno di integrale}
Se una funzione integranda dipende, oltre che dalla variabile di integrazione $\theta$, da un parametro $a$, anche
l'integrale è una funzione di $a$: $F(a) = \int g(a, \theta) \dd\theta$. Per le funzioni regolari che incontreremo, la
derivata dell'integrale è l'integrale della derivata parziale:
\[ \frac{\dd F}{\dd a} = \int \frac{\partial g(a, \theta)}{\partial a} \dd\theta . \]
Il motivo è che l'integrale è un limite di somme (Capitolo~10), e la derivata di una somma è la somma delle derivate.
Per esempio, con $g = (a - \theta)^2$ su $[0, 1]$: $F(a) = a^2 - a + \frac13$, e infatti $F'(a) = 2a - 1 = \int_0^1 2(a
- \theta) \dd\theta$.
\end{richiamo}

\teorema[mmse]{Lo stimatore MMSE è la media condizionata}{
  Con il costo quadratico, lo stimatore bayesiano ottimo è la media a posteriori del parametro:
  \[ \hat\theta_{\text{MMSE}}(\vett{x}) = \int \theta\, f_{\Theta \mid \vett{X}}(\theta \mid \vett{x}) \dd\theta =
     E\left[ \Theta \mid \vett{X} = \vett{x} \right] . \]
}

\begin{dimostrazione}[si annulla la derivata]
\passo{La condizione di minimo} Per il Teorema~\ref{teo:stimatoreottimo} bisogna rendere minima, rispetto a
$\hat\theta(\vett{x})$, la funzione $\int \left( \hat\theta(\vett{x}) - \theta \right)^2 f_{\Theta \mid \vett{X}}(\theta
\mid \vett{x}) \dd\theta$. Si annulla la sua derivata, portandola sotto il segno di integrale:
\[ \frac{\partial}{\partial \hat\theta(\vett{x})} \int \left( \hat\theta(\vett{x}) - \theta \right)^2 f_{\Theta \mid
   \vett{X}}(\theta \mid \vett{x}) \dd\theta = 2 \int \left( \hat\theta(\vett{x}) - \theta \right) f_{\Theta \mid
   \vett{X}}(\theta \mid \vett{x}) \dd\theta = 0 . \]
\passo{Si separano i due termini}
\[ \int \hat\theta(\vett{x})\, f_{\Theta \mid \vett{X}}(\theta \mid \vett{x}) \dd\theta = \int \theta\, f_{\Theta \mid
   \vett{X}}(\theta \mid \vett{x}) \dd\theta . \]
\passo{Si conclude} $\hat\theta(\vett{x})$ non dipende da $\theta$ ed esce dall'integrale, e l'integrale della densità
condizionata vale 1: resta $\hat\theta(\vett{x}) = \int \theta f_{\Theta \mid \vett{X}}(\theta \mid \vett{x})
\dd\theta = E[\Theta \mid \vett{X} = \vett{x}]$, la media condizionata di $\Theta$ all'evento $\{\vett{X} =
\vett{x}\}$.
\passo{È proprio un minimo} Sviluppando il quadrato, la funzione da minimizzare è $\hat\theta^2 - 2\hat\theta\,
E[\Theta \mid \vett{x}] + E[\Theta^2 \mid \vett{x}]$: una parabola in $\hat\theta$ con la concavità verso l'alto (il
costo quadratico è convesso), il cui unico punto a derivata nulla è il vertice, cioè il minimo.
\end{dimostrazione}

% ============================================================
\section{Costo «tutto o niente»: lo stimatore MAP}
% ============================================================

Si consideri ora un costo che non distingue fra errori grandi e piccoli: un errore è gratis se è minore di una
tolleranza piccolissima, e costa 1 altrimenti. Con la funzione rettangolo del Capitolo~11,
\[ C\left( \hat\Theta(\vett{X}) - \Theta \right) = \overline{\Pi}\left( \frac{\hat\Theta(\vett{X}) - \Theta}{\epsilon}
   \right) = \begin{cases} 0 & \abs*{\hat\Theta(\vett{X}) - \Theta} < \frac{\epsilon}{2} , \\ 1 & \text{altrove,}
   \end{cases} \]
dove $\overline{\Pi} = 1 - \mathrm{rect}$ è il complemento del rettangolo: una «tacca» larga $\epsilon$ attorno allo
zero.

\begin{figure}[H]
\centering
\begin{tikzpicture}
\begin{groupplot}[group style={group size=2 by 1, horizontal sep=1.6cm}, width=7cm, height=4.4cm, quaderno,
  xmin=-2.2, xmax=2.3, ymin=0, ymax=1.5, xlabel={$\hat\theta - \theta$}, xtick={-1,1}, ytick={1}]
\nextgroupplot[title={costo quadratico (MMSE)}]
  \addplot[mspiA, domain=-1.2:1.2, samples=60] {x^2};
\nextgroupplot[title={costo a tacca (MAP)}, xtick={-0.3,0.3}, xticklabels={$-\frac{\epsilon}{2}$,$\frac{\epsilon}{2}$}]
  \addplot[mspiB, line width=1.1pt] coordinates {(-2.1,1) (-0.3,1)};
  \addplot[mspiB, line width=1.1pt] coordinates {(-0.3,0) (0.3,0)};
  \addplot[mspiB, line width=1.1pt] coordinates {(0.3,1) (2.1,1)};
  \draw[mspiB, densely dotted] (axis cs:-0.3,0) -- (axis cs:-0.3,1) (axis cs:0.3,0) -- (axis cs:0.3,1);
\end{groupplot}
\end{tikzpicture}
\caption{Le due funzioni di costo dell'errore di stima. Quella quadratica punisce molto gli errori grandi; quella a
tacca considera giusta una stima entro $\frac{\epsilon}{2}$ da $\theta$ e sbagliata ogni altra, comunque lontana.}
\end{figure}

\teorema[map]{Lo stimatore MAP}{
  Con il costo a tacca, per $\epsilon$ arbitrariamente piccolo, lo stimatore bayesiano ottimo è il punto di massimo
  della densità a posteriori:
  \[ \hat\theta_{\text{MAP}}(\vett{x}) = \arg\max_{\theta} f_{\Theta \mid \vett{X}}(\theta \mid \vett{x}) . \]
  Si chiama \term{stima del massimo a posteriori} (MAP).
}

\begin{dimostrazione}
\passo{Il costo medio a posteriori} Il costo vale 1 dappertutto tranne che nella tacca, dove vale 0; quindi il suo
integrale contro la densità a posteriori è 1 meno l'area della densità sulla tacca:
\[ \int C\left( \hat\theta - \theta \right) f_{\Theta \mid \vett{X}}(\theta \mid \vett{x}) \dd\theta = 1 - \int_{\hat\theta
   - \epsilon/2}^{\hat\theta + \epsilon/2} f_{\Theta \mid \vett{X}}(\theta \mid \vett{x}) \dd\theta \simeq 1 -
   \epsilon\, f_{\Theta \mid \vett{X}}(\hat\theta \mid \vett{x}) , \]
perché per $\epsilon$ piccolissimo l'area su un intervallino è base per altezza (Capitolo~10).
\passo{Il minimo} Il costo è minimo quando $f_{\Theta \mid \vett{X}}(\hat\theta \mid \vett{x})$ è massima: si sceglie
$\hat\theta$ nel punto più alto della densità a posteriori.
\end{dimostrazione}

Il rischio bayesiano corrispondente è
\[ R = \sum_{\vett{x} \in \mathcal{X}^n} p_{\vett{X}}(\vett{x}) \int C\left( \hat\theta(\vett{x}) - \theta \right)
   f_{\Theta \mid \vett{X}}(\theta \mid \vett{x}) \dd\theta \qquad\text{o, nel caso continuo,}\qquad R = \iint
   f_{\vett{X}}(\vett{x})\, C\left( \hat\theta(\vett{x}) - \theta \right) f_{\Theta \mid \vett{X}}(\theta \mid \vett{x})
   \dd\theta \dd\vett{x} . \]
I due stimatori guardano la stessa densità a posteriori in due modi diversi: l'MMSE ne prende il \textbf{baricentro}
(la media), il MAP il \textbf{picco} (la moda). Se la densità a posteriori è simmetrica e con un solo picco, come una
gaussiana, i due coincidono.

% ============================================================
\section{Esempio: la moneta misteriosa}
% ============================================================

Una moneta dà testa con una probabilità $\beta$ che non conosciamo affatto: la modelliamo come la realizzazione di una
variabile aleatoria $B \sim \mathcal{U}(0, 1)$, uniforme (nessun valore è a priori più plausibile di un altro). Si lancia
la moneta $n$ volte: i risultati $\vett{X} \in \{0, 1\}^n$ (1 per testa) sono, \emph{sapendo} che $B = \beta$, variabili
di Bernoulli i.i.d. con parametro $\beta$. È il modello \term{Bernoulli composito}. Come nel Capitolo~19, con il peso
di Hamming $w(\vett{x})$ (il numero di «1»):
\[ p_{\vett{X} \mid B}(\vett{x} \mid \beta) = \beta^{w(\vett{x})}\, (1 - \beta)^{n - w(\vett{x})} . \]
Per trovare la densità a posteriori serve l'integrale di questa espressione in $\beta$.

\begin{richiamo}{l'integrale di Eulero e la funzione Gamma}
Per $a$ e $b$ interi non negativi vale
\[ \int_0^1 t^{a} (1 - t)^{b} \dd t = \frac{a!\, b!}{(a + b + 1)!} . \]
\textbf{Perché.} Si chiama $I(a, b)$ l'integrale e lo si calcola per parti (Capitolo~11), derivando $(1 - t)^b$ e
integrando $t^a$:
\[ I(a, b) = \left[ (1 - t)^b \frac{t^{a + 1}}{a + 1} \right]_0^1 + \frac{b}{a + 1} \int_0^1 t^{a + 1} (1 - t)^{b - 1}
   \dd t = \frac{b}{a + 1}\, I(a + 1, b - 1) , \]
perché il termine fra parentesi quadre vale 0 in entrambi gli estremi (se $b \geq 1$). Ogni passo sposta un'unità da
$b$ ad $a$; dopo $b$ passi
\[ I(a, b) = \frac{b}{a + 1} \cdot \frac{b - 1}{a + 2} \cdots \frac{1}{a + b}\; I(a + b, 0) = \frac{b!\, a!}{(a + b)!}
   \cdot \frac{1}{a + b + 1} = \frac{a!\, b!}{(a + b + 1)!} , \]
dove si è usato $I(a + b, 0) = \int_0^1 t^{a+b} \dd t = \frac{1}{a + b + 1}$. Per esempio $\int_0^1 t (1 - t) \dd t =
\frac12 - \frac13 = \frac16 = \frac{1!\, 1!}{3!}$.

\textbf{La funzione Gamma.} Le slide scrivono la stessa formula con la \term{funzione Gamma di Eulero} $\Gamma(x)$,
che estende il fattoriale ai numeri non interi: per gli interi $\Gamma(m) = (m - 1)!$, e in generale $\Gamma(x + 1) =
x\, \Gamma(x)$ (come $m! = m \cdot (m - 1)!$). La \term{funzione Beta} di Eulero è
\[ B(m, k) = \int_0^1 t^{m - 1} (1 - t)^{k - 1} \dd t = \frac{\Gamma(m)\, \Gamma(k)}{\Gamma(m + k)} , \]
che con $m = a + 1$ e $k = b + 1$ è proprio la formula di sopra.
\end{richiamo}

\teorema[bernoullicomposito]{Le stime della moneta misteriosa}{
  Con $B \sim \mathcal{U}(0, 1)$ e $w = w(\vett{x})$ teste in $n$ lanci:
  \begin{enumerate}[label=\alph*)]
    \item la densità a posteriori è $f_{B \mid \vett{X}}(\beta \mid \vett{x}) = \dfrac{(n + 1)!}{w!\, (n - w)!}\,
          \beta^{w} (1 - \beta)^{n - w}$, per $0 \leq \beta \leq 1$;
    \item la stima MMSE è $\hat\beta_{\text{MMSE}}(\vett{x}) = \dfrac{w + 1}{n + 2}$;
    \item la stima MAP è $\hat\beta_{\text{MAP}}(\vett{x}) = \dfrac{w}{n}$.
  \end{enumerate}
}

\begin{dimostrazione}
\passo{La legge dei dati} Mediando sulla variabile $B$, che ha densità 1 su $[0, 1]$, e usando il richiamo con $a = w$
e $b = n - w$:
\[ p_{\vett{X}}(\vett{x}) = \int_0^1 \beta^{w} (1 - \beta)^{n - w}\, f_B(\beta) \dd\beta = \frac{w!\, (n - w)!}{(n +
   1)!} = \frac{\Gamma(w + 1)\, \Gamma(n - w + 1)}{\Gamma(n + 2)} . \]
\passo{a) La densità a posteriori} Per la legge di Bayes, $f_{B \mid \vett{X}}(\beta \mid \vett{x}) =
\frac{p_{\vett{X} \mid B}(\vett{x} \mid \beta)\, f_B(\beta)}{p_{\vett{X}}(\vett{x})}$, con $f_B = 1$.
\passo{b) La media a posteriori} Per il Teorema~\ref{teo:mmse}, usando di nuovo il richiamo, ora con $a = w + 1$:
\[ \hat\beta_{\text{MMSE}}(\vett{x}) = \int_0^1 \beta\, f_{B \mid \vett{X}}(\beta \mid \vett{x}) \dd\beta =
   \frac{(n + 1)!}{w!\, (n - w)!} \int_0^1 \beta^{w + 1} (1 - \beta)^{n - w} \dd\beta = \frac{(n + 1)!}{w!\, (n - w)!}
   \cdot \frac{(w + 1)!\, (n - w)!}{(n + 2)!} = \frac{w + 1}{n + 2} , \]
perché $\frac{(w + 1)!}{w!} = w + 1$ e $\frac{(n + 1)!}{(n + 2)!} = \frac{1}{n + 2}$ (sono le proprietà
$\Gamma(x + 1) = x\,\Gamma(x)$ usate nelle slide).
\passo{c) Il massimo a posteriori} Il fattore costante $\frac{(n + 1)!}{w!(n - w)!}$ è positivo e non dipende da
$\beta$: non sposta il punto di massimo, che è lo stesso di $\beta^{w} (1 - \beta)^{n - w}$, e anche del suo logaritmo
(richiamo su $\arg\max$). Si annulla la derivata:
\[ \frac{\dd}{\dd\beta} \left[ w \ln \beta + (n - w) \ln (1 - \beta) \right] = \frac{w}{\beta} - \frac{n - w}{1 - \beta} =
   0 \quad\iff\quad (1 - \beta)\, w = \beta\, (n - w) \quad\iff\quad w = \beta n , \]
nell'ipotesi $\beta \neq 0, 1$ (se $w = 0$ o $w = n$ il massimo è in un estremo, e vale ancora $\beta = \frac{w}{n}$).
\end{dimostrazione}

La stima MAP è la frequenza delle teste: la risposta più spontanea. La stima MMSE è leggermente diversa: è come se ai
lanci veri si aggiungessero una testa e una croce immaginarie, che rappresentano la densità a priori uniforme (centrata
in $\frac12$). Con pochi dati la differenza conta: dopo 3 teste in 3 lanci il MAP dice $\hat\beta = 1$, cioè «la moneta
dà sempre testa», mentre l'MMSE dice, più prudente, $\frac{4}{5}$. È la \term{regola di successione di Laplace}, che il
matematico usò per il famoso problema dell'alba: avendo visto sorgere il sole $n$ giorni di fila, la probabilità che
sorga domani è $\frac{n + 1}{n + 2}$.

\begin{figure}[H]
\centering
\begin{tikzpicture}
\begin{axis}[width=14cm, height=6cm, xmin=0, xmax=1, ymin=0, ymax=6.8, xtick={0,0.2,0.4,0.5,0.6,0.8,1},
  xticklabels={$0$,$0{,}2$,$0{,}4$,$0{,}5$,$0{,}6$,$0{,}8$,$1$}, ytick={2,4,6}, xlabel=$\beta$,
  ylabel={$f_{B \mid \vett{X}}(\beta \mid \vett{x})$}, axis lines=left, tick label style={font=\scriptsize},
  label style={font=\small}, grid=major, grid style={mspigray!20},
  legend style={font=\scriptsize, at={(0.02,0.98)}, anchor=north west}]
  \addplot[mspigray, line width=1pt] coordinates {(0,1) (1,1)};
  \addplot[mspiC, line width=1.1pt] coordinates {(0.000,0.0000) (0.005,0.0000) (0.010,0.0000) (0.015,0.0000) (0.020,0.0000) (0.025,0.0001) (0.030,0.0001) (0.035,0.0002) (0.040,0.0003) (0.045,0.0004) (0.050,0.0005) (0.055,0.0007) (0.060,0.0009) (0.065,0.0011) (0.070,0.0014) (0.075,0.0017) (0.080,0.0020) (0.085,0.0025) (0.090,0.0029) (0.095,0.0034) (0.100,0.0040) (0.105,0.0046) (0.110,0.0053) (0.115,0.0061) (0.120,0.0069) (0.125,0.0078) (0.130,0.0088) (0.135,0.0098) (0.140,0.0110) (0.145,0.0122) (0.150,0.0135) (0.155,0.0149) (0.160,0.0164) (0.165,0.0180) (0.170,0.0197) (0.175,0.0214) (0.180,0.0233) (0.185,0.0253) (0.190,0.0274) (0.195,0.0297) (0.200,0.0320) (0.205,0.0345) (0.210,0.0370) (0.215,0.0398) (0.220,0.0426) (0.225,0.0456) (0.230,0.0487) (0.235,0.0519) (0.240,0.0553) (0.245,0.0588) (0.250,0.0625) (0.255,0.0663) (0.260,0.0703) (0.265,0.0744) (0.270,0.0787) (0.275,0.0832) (0.280,0.0878) (0.285,0.0926) (0.290,0.0976) (0.295,0.1027) (0.300,0.1080) (0.305,0.1135) (0.310,0.1192) (0.315,0.1250) (0.320,0.1311) (0.325,0.1373) (0.330,0.1437) (0.335,0.1504) (0.340,0.1572) (0.345,0.1643) (0.350,0.1715) (0.355,0.1790) (0.360,0.1866) (0.365,0.1945) (0.370,0.2026) (0.375,0.2109) (0.380,0.2195) (0.385,0.2283) (0.390,0.2373) (0.395,0.2465) (0.400,0.2560) (0.405,0.2657) (0.410,0.2757) (0.415,0.2859) (0.420,0.2964) (0.425,0.3071) (0.430,0.3180) (0.435,0.3293) (0.440,0.3407) (0.445,0.3525) (0.450,0.3645) (0.455,0.3768) (0.460,0.3893) (0.465,0.4022) (0.470,0.4153) (0.475,0.4287) (0.480,0.4424) (0.485,0.4563) (0.490,0.4706) (0.495,0.4851) (0.500,0.5000) (0.505,0.5152) (0.510,0.5306) (0.515,0.5464) (0.520,0.5624) (0.525,0.5788) (0.530,0.5955) (0.535,0.6125) (0.540,0.6299) (0.545,0.6475) (0.550,0.6655) (0.555,0.6838) (0.560,0.7025) (0.565,0.7214) (0.570,0.7408) (0.575,0.7604) (0.580,0.7804) (0.585,0.8008) (0.590,0.8215) (0.595,0.8426) (0.600,0.8640) (0.605,0.8858) (0.610,0.9079) (0.615,0.9304) (0.620,0.9533) (0.625,0.9766) (0.630,1.0002) (0.635,1.0242) (0.640,1.0486) (0.645,1.0733) (0.650,1.0985) (0.655,1.1240) (0.660,1.1500) (0.665,1.1763) (0.670,1.2031) (0.675,1.2302) (0.680,1.2577) (0.685,1.2857) (0.690,1.3140) (0.695,1.3428) (0.700,1.3720) (0.705,1.4016) (0.710,1.4316) (0.715,1.4621) (0.720,1.4930) (0.725,1.5243) (0.730,1.5561) (0.735,1.5883) (0.740,1.6209) (0.745,1.6540) (0.750,1.6875) (0.755,1.7215) (0.760,1.7559) (0.765,1.7908) (0.770,1.8261) (0.775,1.8619) (0.780,1.8982) (0.785,1.9349) (0.790,1.9722) (0.795,2.0098) (0.800,2.0480) (0.805,2.0866) (0.810,2.1258) (0.815,2.1654) (0.820,2.2055) (0.825,2.2461) (0.830,2.2871) (0.835,2.3287) (0.840,2.3708) (0.845,2.4134) (0.850,2.4565) (0.855,2.5001) (0.860,2.5442) (0.865,2.5889) (0.870,2.6340) (0.875,2.6797) (0.880,2.7259) (0.885,2.7726) (0.890,2.8199) (0.895,2.8677) (0.900,2.9160) (0.905,2.9649) (0.910,3.0143) (0.915,3.0642) (0.920,3.1148) (0.925,3.1658) (0.930,3.2174) (0.935,3.2696) (0.940,3.3223) (0.945,3.3756) (0.950,3.4295) (0.955,3.4839) (0.960,3.5389) (0.965,3.5945) (0.970,3.6507) (0.975,3.7074) (0.980,3.7648) (0.985,3.8227) (0.990,3.8812) (0.995,3.9403) (1.000,4.0000)};
  \addplot[mspiA, line width=1.1pt] coordinates {(0.000,0.0000) (0.005,0.0000) (0.010,0.0000) (0.015,0.0000) (0.020,0.0000) (0.025,0.0000) (0.030,0.0000) (0.035,0.0000) (0.040,0.0000) (0.045,0.0000) (0.050,0.0000) (0.055,0.0000) (0.060,0.0000) (0.065,0.0000) (0.070,0.0000) (0.075,0.0000) (0.080,0.0000) (0.085,0.0000) (0.090,0.0000) (0.095,0.0001) (0.100,0.0001) (0.105,0.0001) (0.110,0.0002) (0.115,0.0002) (0.120,0.0003) (0.125,0.0004) (0.130,0.0005) (0.135,0.0007) (0.140,0.0009) (0.145,0.0011) (0.150,0.0014) (0.155,0.0017) (0.160,0.0021) (0.165,0.0026) (0.170,0.0031) (0.175,0.0037) (0.180,0.0045) (0.185,0.0053) (0.190,0.0063) (0.195,0.0074) (0.200,0.0087) (0.205,0.0101) (0.210,0.0117) (0.215,0.0136) (0.220,0.0156) (0.225,0.0179) (0.230,0.0205) (0.235,0.0234) (0.240,0.0266) (0.245,0.0301) (0.250,0.0340) (0.255,0.0383) (0.260,0.0430) (0.265,0.0481) (0.270,0.0537) (0.275,0.0598) (0.280,0.0665) (0.285,0.0737) (0.290,0.0815) (0.295,0.0899) (0.300,0.0990) (0.305,0.1088) (0.310,0.1193) (0.315,0.1306) (0.320,0.1426) (0.325,0.1555) (0.330,0.1692) (0.335,0.1838) (0.340,0.1993) (0.345,0.2158) (0.350,0.2332) (0.355,0.2517) (0.360,0.2712) (0.365,0.2917) (0.370,0.3133) (0.375,0.3361) (0.380,0.3599) (0.385,0.3850) (0.390,0.4112) (0.395,0.4386) (0.400,0.4671) (0.405,0.4969) (0.410,0.5280) (0.415,0.5602) (0.420,0.5937) (0.425,0.6285) (0.430,0.6645) (0.435,0.7017) (0.440,0.7401) (0.445,0.7798) (0.450,0.8206) (0.455,0.8627) (0.460,0.9059) (0.465,0.9502) (0.470,0.9956) (0.475,1.0421) (0.480,1.0896) (0.485,1.1381) (0.490,1.1876) (0.495,1.2379) (0.500,1.2891) (0.505,1.3410) (0.510,1.3936) (0.515,1.4469) (0.520,1.5008) (0.525,1.5551) (0.530,1.6099) (0.535,1.6650) (0.540,1.7203) (0.545,1.7758) (0.550,1.8313) (0.555,1.8867) (0.560,1.9420) (0.565,1.9970) (0.570,2.0517) (0.575,2.1058) (0.580,2.1593) (0.585,2.2121) (0.590,2.2641) (0.595,2.3150) (0.600,2.3649) (0.605,2.4135) (0.610,2.4608) (0.615,2.5066) (0.620,2.5507) (0.625,2.5932) (0.630,2.6337) (0.635,2.6722) (0.640,2.7086) (0.645,2.7427) (0.650,2.7744) (0.655,2.8036) (0.660,2.8302) (0.665,2.8540) (0.670,2.8750) (0.675,2.8930) (0.680,2.9079) (0.685,2.9197) (0.690,2.9282) (0.695,2.9334) (0.700,2.9351) (0.705,2.9333) (0.710,2.9280) (0.715,2.9191) (0.720,2.9065) (0.725,2.8903) (0.730,2.8703) (0.735,2.8466) (0.740,2.8191) (0.745,2.7880) (0.750,2.7531) (0.755,2.7146) (0.760,2.6724) (0.765,2.6267) (0.770,2.5775) (0.775,2.5248) (0.780,2.4689) (0.785,2.4098) (0.790,2.3476) (0.795,2.2825) (0.800,2.2146) (0.805,2.1441) (0.810,2.0712) (0.815,1.9961) (0.820,1.9191) (0.825,1.8402) (0.830,1.7598) (0.835,1.6782) (0.840,1.5955) (0.845,1.5120) (0.850,1.4282) (0.855,1.3441) (0.860,1.2602) (0.865,1.1768) (0.870,1.0941) (0.875,1.0124) (0.880,0.9322) (0.885,0.8536) (0.890,0.7771) (0.895,0.7029) (0.900,0.6314) (0.905,0.5627) (0.910,0.4973) (0.915,0.4353) (0.920,0.3770) (0.925,0.3227) (0.930,0.2724) (0.935,0.2265) (0.940,0.1849) (0.945,0.1478) (0.950,0.1152) (0.955,0.0871) (0.960,0.0635) (0.965,0.0441) (0.970,0.0288) (0.975,0.0173) (0.980,0.0092) (0.985,0.0040) (0.990,0.0012) (0.995,0.0002) (1.000,0.0000)};
  \addplot[mspiB, line width=1.1pt] coordinates {(0.000,0.0000) (0.005,0.0000) (0.010,0.0000) (0.015,0.0000) (0.020,0.0000) (0.025,0.0000) (0.030,0.0000) (0.035,0.0000) (0.040,0.0000) (0.045,0.0000) (0.050,0.0000) (0.055,0.0000) (0.060,0.0000) (0.065,0.0000) (0.070,0.0000) (0.075,0.0000) (0.080,0.0000) (0.085,0.0000) (0.090,0.0000) (0.095,0.0000) (0.100,0.0000) (0.105,0.0000) (0.110,0.0000) (0.115,0.0000) (0.120,0.0000) (0.125,0.0000) (0.130,0.0000) (0.135,0.0000) (0.140,0.0000) (0.145,0.0000) (0.150,0.0000) (0.155,0.0000) (0.160,0.0000) (0.165,0.0000) (0.170,0.0000) (0.175,0.0000) (0.180,0.0000) (0.185,0.0000) (0.190,0.0000) (0.195,0.0000) (0.200,0.0000) (0.205,0.0000) (0.210,0.0000) (0.215,0.0000) (0.220,0.0000) (0.225,0.0000) (0.230,0.0000) (0.235,0.0000) (0.240,0.0000) (0.245,0.0000) (0.250,0.0000) (0.255,0.0000) (0.260,0.0000) (0.265,0.0000) (0.270,0.0000) (0.275,0.0000) (0.280,0.0000) (0.285,0.0000) (0.290,0.0000) (0.295,0.0000) (0.300,0.0000) (0.305,0.0000) (0.310,0.0000) (0.315,0.0000) (0.320,0.0000) (0.325,0.0000) (0.330,0.0000) (0.335,0.0000) (0.340,0.0000) (0.345,0.0000) (0.350,0.0000) (0.355,0.0000) (0.360,0.0000) (0.365,0.0001) (0.370,0.0001) (0.375,0.0001) (0.380,0.0002) (0.385,0.0002) (0.390,0.0003) (0.395,0.0005) (0.400,0.0006) (0.405,0.0009) (0.410,0.0012) (0.415,0.0016) (0.420,0.0021) (0.425,0.0028) (0.430,0.0037) (0.435,0.0049) (0.440,0.0064) (0.445,0.0083) (0.450,0.0107) (0.455,0.0137) (0.460,0.0175) (0.465,0.0222) (0.470,0.0280) (0.475,0.0352) (0.480,0.0440) (0.485,0.0547) (0.490,0.0677) (0.495,0.0833) (0.500,0.1020) (0.505,0.1242) (0.510,0.1506) (0.515,0.1817) (0.520,0.2181) (0.525,0.2606) (0.530,0.3098) (0.535,0.3665) (0.540,0.4316) (0.545,0.5058) (0.550,0.5899) (0.555,0.6848) (0.560,0.7912) (0.565,0.9098) (0.570,1.0413) (0.575,1.1861) (0.580,1.3447) (0.585,1.5174) (0.590,1.7041) (0.595,1.9047) (0.600,2.1189) (0.605,2.3459) (0.610,2.5848) (0.615,2.8343) (0.620,3.0930) (0.625,3.3588) (0.630,3.6296) (0.635,3.9029) (0.640,4.1759) (0.645,4.4456) (0.650,4.7087) (0.655,4.9619) (0.660,5.2016) (0.665,5.4243) (0.670,5.6265) (0.675,5.8049) (0.680,5.9563) (0.685,6.0777) (0.690,6.1668) (0.695,6.2213) (0.700,6.2397) (0.705,6.2210) (0.710,6.1649) (0.715,6.0715) (0.720,5.9418) (0.725,5.7774) (0.730,5.5805) (0.735,5.3537) (0.740,5.1006) (0.745,4.8249) (0.750,4.5306) (0.755,4.2223) (0.760,3.9044) (0.765,3.5816) (0.770,3.2584) (0.775,2.9391) (0.780,2.6276) (0.785,2.3277) (0.790,2.0424) (0.795,1.7745) (0.800,1.5259) (0.805,1.2980) (0.810,1.0919) (0.815,0.9078) (0.820,0.7456) (0.825,0.6045) (0.830,0.4835) (0.835,0.3812) (0.840,0.2961) (0.845,0.2264) (0.850,0.1702) (0.855,0.1257) (0.860,0.0910) (0.865,0.0646) (0.870,0.0449) (0.875,0.0305) (0.880,0.0202) (0.885,0.0130) (0.890,0.0081) (0.895,0.0049) (0.900,0.0029) (0.905,0.0016) (0.910,0.0009) (0.915,0.0004) (0.920,0.0002) (0.925,0.0001) (0.930,0.0000) (0.935,0.0000) (0.940,0.0000) (0.945,0.0000) (0.950,0.0000) (0.955,0.0000) (0.960,0.0000) (0.965,0.0000) (0.970,0.0000) (0.975,0.0000) (0.980,0.0000) (0.985,0.0000) (0.990,0.0000) (0.995,0.0000) (1.000,0.0000)};
  \legend{{a priori, $\mathcal{U}(0, 1)$}, {$n = 3$, $w = 3$}, {$n = 10$, $w = 7$}, {$n = 50$, $w = 35$}}
  \draw[mspiA, dashed] (axis cs:0.6667,0) -- (axis cs:0.6667,2.52);
  \draw[mspiA, densely dotted, line width=0.9pt] (axis cs:0.7,0) -- (axis cs:0.7,2.935);
  \draw[-{Stealth}, mspiA!80!black] (axis cs:0.45,4.3) node[above, font=\scriptsize] {stima MMSE, $\frac{8}{12}$}
    -- (axis cs:0.655,2.6);
  \draw[-{Stealth}, mspiA!80!black] (axis cs:0.9,4.3) node[above, font=\scriptsize] {stima MAP, $0{,}7$}
    -- (axis cs:0.712,3.0);
\end{axis}
\end{tikzpicture}
\caption{Le densità a posteriori della probabilità di testa. Partendo dalla densità uniforme, i dati la spostano e la
stringono: con 35 teste su 50 lanci il parametro è quasi certamente fra $0{,}55$ e $0{,}85$. Per $n = 10$ e $w = 7$ sono
segnate la media (stima MMSE, $\frac{8}{12} \approx 0{,}667$) e il picco (stima MAP, $0{,}7$).}
\end{figure}

% ============================================================
\section{Le prestazioni dei due stimatori}
% ============================================================

Quale dei due stimatori sbaglia meno, in media? Calcoliamo media ed errore quadratico medio di entrambi, mediando sia
sui dati sia sul parametro. Servono alcuni momenti, che si calcolano tutti con il teorema della media condizionata:
prima si media sapendo $B = \beta$ (e allora $w(\vett{X})$ è binomiale con parametri $n$ e $\beta$), poi su $B$.
\begin{itemize}
  \item Per la variabile uniforme: $E[B] = \int_0^1 \beta \dd\beta = \frac12$, $E[B^2] = \int_0^1 \beta^2 \dd\beta =
        \left[ \frac{\beta^3}{3} \right]_0^1 = \frac13$, e $\sigma_B^2 = \frac13 - \frac14 = \frac{1}{12}$.
  \item La media del peso: $E[w(\vett{X})] = E_B\left[ E[w(\vett{X}) \mid B] \right] = E_B[nB] = \frac{n}{2}$.
  \item Il valore quadratico medio del peso: sapendo $B$, la binomiale ha varianza $nB(1 - B)$ e media $nB$, e il valore
        quadratico medio è varianza più quadrato della media:
        \[ E\left[ w^2(\vett{X}) \right] = E_B\left[ nB(1 - B) + n^2 B^2 \right] = n \left( \frac12 - \frac13 \right) +
           \frac{n^2}{3} = \frac{n}{6} + \frac{n^2}{3} , \]
        e quindi la varianza (non condizionata) del peso è $\sigma_{w}^2 = \frac{n}{6} + \frac{n^2}{3} - \frac{n^2}{4} =
        \frac{n}{6} + \frac{n^2}{12}$.
  \item Il prodotto: $E[w(\vett{X})\, B] = E_B\left[ B\, E[w(\vett{X}) \mid B] \right] = E_B[n B^2] = \frac{n}{3}$.
\end{itemize}

\teorema[prestazioni]{Media ed errore quadratico medio dei due stimatori}{
  Entrambi gli stimatori hanno media $\frac12$, come il parametro: $E\left[ \hat\beta_{\text{MMSE}}(\vett{X}) \right] =
  E\left[ \hat\beta_{\text{MAP}}(\vett{X}) \right] = \frac12$. I loro errori quadratici medi sono
  \[ E\left[ \left( \hat\beta_{\text{MMSE}}(\vett{X}) - B \right)^2 \right] = \frac{1}{6(n + 2)} , \qquad E\left[
     \left( \hat\beta_{\text{MAP}}(\vett{X}) - B \right)^2 \right] = \frac{1}{6n} . \]
}

\begin{dimostrazione}
\passo{Le medie} Per la linearità: $E\left[ \frac{w + 1}{n + 2} \right] = \frac{n/2 + 1}{n + 2} = \frac{(n + 2)/2}{n +
2} = \frac12$ e $E\left[ \frac{w}{n} \right] = \frac{n/2}{n} = \frac12$.
\passo{L'errore quadratico in tre pezzi} Per entrambi gli stimatori, sviluppando il quadrato,
$E[(\hat\beta - B)^2] = E[\hat\beta^2] + E[B^2] - 2E[\hat\beta B]$, con $E[B^2] = \frac13$.
\passo{MMSE} Con i momenti calcolati sopra:
\[ E\left[ \hat\beta_{\text{MMSE}}^2 \right] = \frac{E[w^2] + 2E[w] + 1}{(n + 2)^2} = \frac{\frac{n}{6} + \frac{n^2}{3}
   + n + 1}{(n + 2)^2} = \frac{2n^2 + 7n + 6}{6(n + 2)^2} = \frac{(2n + 3)(n + 2)}{6(n + 2)^2} = \frac{2n + 3}{6(n +
   2)} , \]
\[ E\left[ \hat\beta_{\text{MMSE}}\, B \right] = \frac{E[wB] + E[B]}{n + 2} = \frac{\frac{n}{3} + \frac12}{n + 2} =
   \frac{2n + 3}{6(n + 2)} . \]
Quindi l'errore è $\frac{2n + 3}{6(n + 2)} + \frac13 - 2 \cdot \frac{2n + 3}{6(n + 2)} = \frac13 - \frac{2n + 3}{6(n +
2)} = \frac{2(n + 2) - (2n + 3)}{6(n + 2)} = \frac{1}{6(n + 2)}$.
\passo{MAP}
\[ E\left[ \hat\beta_{\text{MAP}}^2 \right] = \frac{E[w^2]}{n^2} = \frac{1}{6n} + \frac13 , \qquad E\left[
   \hat\beta_{\text{MAP}}\, B \right] = \frac{E[wB]}{n} = \frac13 , \]
e l'errore è $\frac{1}{6n} + \frac13 + \frac13 - \frac23 = \frac{1}{6n}$.
\end{dimostrazione}

Le slide arrivano agli stessi risultati con passaggi più lunghi (slide 64--73), sviluppando $(n + 2)^2$ e i prodotti.
Come c'era da aspettarsi, per ogni valore di $n$ l'errore quadratico medio dello stimatore MMSE è minore di quello del
MAP: è il migliore possibile, per costruzione. La differenza è grande con pochi dati e sparisce quando i dati sono
tanti:
\begin{center}
\renewcommand{\arraystretch}{1.3}
\begin{tabular}{@{}c|ccccc@{}}
  $n$ & 1 & 2 & 5 & 10 & 100 \\ \hline
  errore MMSE, $\frac{1}{6(n+2)}$ & $0{,}056$ & $0{,}042$ & $0{,}024$ & $0{,}0139$ & $0{,}00163$ \\
  errore MAP, $\frac{1}{6n}$ & $0{,}167$ & $0{,}083$ & $0{,}033$ & $0{,}0167$ & $0{,}00167$
\end{tabular}
\end{center}

\dfn{Consistenza di uno stimatore}{
  Uno stimatore $\hat B_n(\vett{X})$, calcolato da $n$ dati, si dice
  \begin{itemize}
    \item \term{consistente in media quadratica} (\emph{MS-consistent}) se il suo errore quadratico medio tende a zero
          per $n \to \infty$;
    \item \term{consistente} (in senso debole) se tende al parametro in probabilità: per ogni $\epsilon > 0$,
          $\lim_{n \to \infty} P\left\{ \abs*{\hat B_n(\vett{X}) - B} > \epsilon \right\} = 0$.
  \end{itemize}
}

Entrambi gli stimatori della moneta sono consistenti in media quadratica, perché $\frac{1}{6(n+2)}$ e $\frac{1}{6n}$
tendono a zero. E quindi sono anche consistenti in senso debole: per la disuguaglianza di Markov, come nel
Capitolo~18,
\[ P\left\{ \abs*{\hat B_n(\vett{X}) - B} > \epsilon \right\} = P\left\{ \left( \hat B_n(\vett{X}) - B \right)^2 >
   \epsilon^2 \right\} \leq \frac{E\left[ \left( \hat B_n(\vett{X}) - B \right)^2 \right]}{\epsilon^2} =
   \frac{\mathrm{MSE}\left( \hat B_n \right)}{\epsilon^2} \;\xrightarrow{\;n \to \infty\;}\; 0 , \]
dove $\mathrm{MSE}$ (\emph{Mean Square Error}) indica l'errore quadratico medio: al crescere di $n$ il secondo membro
tende a zero, e il primo, che è una probabilità, non è negativo.

\gbox{In sintesi}{mspiA!80!black}{
\begin{itemize}
  \item \textbf{Stima bayesiana}: il parametro è una variabile $\Theta$ con densità a priori $f_\Theta$; i dati la
        trasformano nella densità a posteriori $f_{\Theta \mid \vett{X}}$ (legge di Bayes).
  \item Lo \textbf{stimatore ottimo} minimizza il costo medio a posteriori, per ogni $\vett{x}$:
        $\hat\theta = \arg\min \int C(\hat\theta - \theta) f_{\Theta \mid \vett{X}}(\theta \mid \vett{x}) \dd\theta$.
  \item Costo quadratico: \textbf{MMSE} $= E[\Theta \mid \vett{X} = \vett{x}]$, il baricentro della densità a posteriori.
        Costo a tacca: \textbf{MAP} $= \arg\max f_{\Theta \mid \vett{X}}(\theta \mid \vett{x})$, il suo picco.
  \item Moneta con $B \sim \mathcal{U}(0,1)$ e $w$ teste in $n$ lanci: $\hat\beta_{\text{MMSE}} = \frac{w + 1}{n + 2}$,
        $\hat\beta_{\text{MAP}} = \frac{w}{n}$; errori quadratici $\frac{1}{6(n+2)} < \frac{1}{6n}$, entrambi $\to 0$:
        stimatori consistenti.
  \item $\int_0^1 t^a (1 - t)^b \dd t = \frac{a!\, b!}{(a + b + 1)!}$ (funzioni Beta e Gamma di Eulero).
\end{itemize}
}
